\subsection{The Redistricting Problem}

Congressional redistricting has become synonymous with partisan manipulation. Following \emph{Baker v. Carr} (1962) and \emph{Reynolds v. Sims} (1964), the Supreme Court established the ``one person, one vote'' doctrine, requiring districts to be population-balanced \citep{baker1962carr, reynolds1964sims}. While this achieved greater equality than pre-1960s malapport ionment, it created a new problem: if districts must be redrawn every decade to maintain population equality, who controls the pen controls political power \citep{stephanopoulos2015partisan}.

The efficiency gap literature demonstrates how even  population-balanced districts can be gerrymandered through strategic ``packing'' (concentrating opposition voters) and ``cracking'' (dispersing them across districts) \citep{stephanopoulos2015partisan, mcghee2014measuring}. Recent cases like \emph{Rucho v. Common Cause} (2019) held that partisan gerrymandering claims are non-justiciable, leaving no federal remedy \citep{rucho2019}. State reforms---independent commissions, algorithmic fairness criteria, court challenges---provide partial relief but do not eliminate the fundamental problem: arbitrary boundaries invite manipulation.

\subsection{Natural Boundaries and Communities of Interest}

Political scientists have long argued for respecting ``communities of interest'' in redistricting: areas with shared economic, social, or cultural characteristics that merit unified representation \citep{morrill1981political}. Counties are the quintessential American community of interest. They predate congressional districts, with most county boundaries established in the 18th or 19th century. Counties provide unified governance (sheriffs, courts, health services), economic integration (labor markets, commerce), and civic identity (residents identify as ``from Cook County'' or ``from Orange County'') \citep{benton2002counties}.

Yet redistricting systematically fragments counties. \citet{altman2005contiguity} found that only 37\% of congressional districts respect county boundaries, with most crossing multiple jurisdictions. This fragmentation dilutes local political identity and creates accountability problems: representatives serve portions of counties rather than coherent communities \citep{butler2016local}.

\subsection{Subsidiarity and Local Autonomy}

The principle of subsidiarity---that decisions should be made at the most local level possible---has deep roots in political theory from \citet{tocqueville1835democracy} to modern federalism scholars \citep{ostrom1999polycentricity}. County governance embodies subsidiarity: local control over education, law enforcement, public health, and land use. Yet congressional representation---ostensibly ``local'' representatives to the national government---is controlled by state legislatures that may fragment counties for state-level political goals.

This creates a mismatch: counties govern themselves but cannot represent themselves at the federal level. A state legislaure in Sacramento or Albany can divide Los Angeles or Brooklyn for strategic purposes unrelated to local community cohesion. County-based federal representation would align governance and representation, respecting local autonomy in the same way the Constitution respects state autonomy in the Senate.

\subsection{Historical Precedents for Natural Boundary Representation}

The U.S. Senate provides the clearest precedent: states, as pre-existing political units, receive representation based on their boundaries rather than population-optimized districts. Though the Senate violates one person, one vote (Wyoming has 1 senator per 288,000 people, California 1 per 20 million), this is constitutionally enshrined as respecting state sovereignty \citep{lee2010sizing}.

European federal systems provide additional models. The German Bundesrat represents L\"ander (states) with weighted voting: larger states cast more votes than smaller ones, but all vote as blocks based on their boundaries \citep{gunlicks2003german}. Switzerland's Council of States represents cantons, though equally rather than proportionally. While these systems differ from direct proportional representation, they demonstrate that \emph{natural boundary-based representation} is a viable alternative to population-optimized districting.

The key distinction is between \emph{geographic units of governance} (states, counties) and \emph{arbitrary lines drawn for electoral convenience} (congressional districts). The former have coherent political identities; the latter exist only on maps.
